\section{Proof status, reproduction, and further questions}
\label{sec:replay}

\subsection{What Lean checks}

\noindent\begin{tabularx}{\linewidth}{@{}p{.23\linewidth}YY@{}}
\toprule
Component & Kernel-checked result & Remaining object derivation\\
\midrule
Finite observations & Object iPRA sparse replay at variable time. Lean window certification, composition of prefix dependence and no-window example. & Reify the field operations for each reader and prove its finite use and extensionality contracts.\\[3pt]
Zero selection & Lean brackets, dyadic widths, cross-stage midpoint bound and residual schedule. Closed iPRA theorem for the three sign branches. & Compile the dyadic transition and run the stage induction at fixed test precision.\\[3pt]
Differentiation & Lean local signed bounds. Closed iPRA theorem deriving both child/parent comparisons from the adjacent-child comparison. & Sum the errors over a level path, derive the cell estimates, apply the tagged FTC and pass to the names.\\[3pt]
Interval laws & Lean finite Abel formula and optimal dual potential. Object iPRA variable-length prefix maximum and its common upper bound. & Code the sparse coupling and verify the two marginal scans and reconstruction contracts.\\[3pt]
Metric partitions and quantiles & Human proofs and explicit programs or finite bounds. Exact-arithmetic regression for the density $1/2+x$. & Reify rounding of finite potentials and compose the ramp accuracy proof with the midpoint invariant.\\
\bottomrule
\end{tabularx}

The archive contains four standalone Lean modules using the Lean 4.32 standard
distribution, and four modules importing ProveIt. In the latter, the exported
types mention \texttt{Prf PRAAx} literally. The zero and slope lemmas close
their seven numerical variables in the object language. The prefix scan is
a derivation constructor at arbitrary terms and hypotheses with the raw sample
body fixed syntactically. Sparse replay has the same schema form for fixed reader instructions. All four use intuitionistic object rules. The
ordinary Lean metatheory used by the upstream proof infrastructure reports
\texttt{propext}, \texttt{Classical.choice}, and \texttt{Quot.sound} for the
audited object results. These constants occur in the host implementation of
proof constructions. The object judgments use the repository's intuitionistic axiom family.

\subsection{Dependency order for completing the formalization}

\begin{enumerate}[leftmargin=*]
\item Define a total encoding of lists, reduced dyadic sites and signed
rationals on top of \texttt{PTm.toPR}, \texttt{recDef} and the pairing and
order modules. Prove normalization preserves equality and cross-multiplication
comparison, addition, multiplication, absolute value and clipping. A unit
interval site may be coded by numerator and power-of-two denominator with a
clipping operation. Record validity checks for arbitrary input codes. Fix input
evaluator symbols and numerical parameters when forming a sequent.
\item Formalize regularity, point equality, forward and inverse moduli, and
value-table inequalities as universally closed formulas with the rational
operations expanded. For instance, the slope contract for $A,b$ quantifies
over the natural triple $(n,j,p)$ and, under $0<j<2^n$, tests the two signed
cross-multiplied versions of \eqref{eq:dd-grid}. Trace verification and legality
of queries use bounded scans. For a fixed reader, substitute a PR symbol for
each transition and apply the checked sparse replay theorem to its two state codes at the requested terminal slot.
\item For the roots, prove the two defining equations of the integer endpoint
recurrence, instantiate the checked three-way sequent and use the width and
sign invariants at fixed outer precision. At a single final precision choose
the starting stage from $\mu$ and finish with \eqref{eq:zero-forward}.
For differentiation use the ancestor-induction and slack bounds beside
\eqref{eq:dd-finite}. Prove the signed cell estimate before using the
repository's tagged-sum derivations. Its two bounds have hypotheses for every
cell, so the passage must include the denominator scale for each finite
tagged term and both endpoints. A product of the finite denominators is a
sufficient common positive scale.
\item For laws, code mergers and residual atom weights by the integer
numerators at a common positive denominator. At each transfer maintain the
two prescribed masses less the emitted row and column sums. Transfer their
minimum and advance at least one cursor. The sum of the remaining numbers of
sites decreases; iteration for $s+t$ steps with an absorbing final state
is PR. At a fixed cut maintain the cumulative crossing flow and use the common
mass-coordinate formula of Lemma~\ref{lem:dual}. Use the repository's
\texttt{Sums.abel\_summation} for the signed equality after denominator
clearing, splitting signed terms into their natural components. The finite
host specification checks the index and boundary conventions.
\item Apply \texttt{prefix\_upper} at length $i+1$ to the repaired ramp readings.
Derive the lower bound using \texttt{covers\_latest}. Prove the prefix
expansion of the interpolant at fixed auxiliary precision as in
Section~\ref{sec:laws}, then instantiate the maximizing polygon constructed
from the two outputs. For a new partition language reify its constructor
identities (in particular the sum of the weights), the support comparison and
finite minimum extension. A dyadic cube uses bounded products of $d$ hats.
For finite dual search enumerate the bounded ordinate array, sort its
fractional jumps, and telescope the positive rational layer lengths. This
supplies the selection proof for the finite maximum. For the quantile
specialize the law evaluator to the right ramp with the displayed level
$2p+d+4$, weaken the input and output estimates to the root contracts,
and apply the signed stage induction from step 3.
\end{enumerate}
These steps specify the inputs to each open lemma, the bounded witness, the
induction matrix or its rational inequality, and the existing module it uses.
The full operators currently have human arithmetic proofs and partially
mechanized components.

\subsection{Sources and further research}

The source inventory accompanying the paper records the pinned repository
paths. The main external references are
A. Pauly, \emph{On the topological aspects of the theory of represented
spaces} (2016), \href{https://arxiv.org/abs/1204.3763}{author preprint};
A. Kawamura and S. Cook, \emph{Complexity theory for operators in analysis},
\href{https://arxiv.org/abs/1305.0453}{author preprint};
M. Hoyrup and C. Rojas, \emph{Computability of probability measures and
Martin-L\"of randomness over metric spaces} (2009),
\href{https://doi.org/10.1016/j.ic.2008.12.009}{published paper}; and
S. S. Vallander, \emph{Calculation of the Wasserstein distance between
probability distributions on the line} (1973, English translation 1974),
\href{https://doi.org/10.1137/1118101}{published paper}.
The one-dimensional transport formula is classical in this literature. Its
use here supplies finite certificates and precision schedules for an
iPRA representation change.

\begin{itemize}
\item On a metric tree, deleting an edge separates the sites into two
components. Can subtree probes replace the line's cumulative ramps and allow
precision to depend on weighted branch length in place of the total number
of cells?
\item Can the dyadic second-difference test use locally varying tail certificates
to construct derivatives on certified subintervals? At a shared endpoint the
two one-sided slope bounds must still give a compatible value.
\item For structured finite laws, which mesh and mass queries can be shared
between successive quantile bits? The worked example uses prefix moments.
A general adaptive version needs a verifiable bound on the cost of each
compressed query and a proof that it realizes the declared finite law.
\end{itemize}
